%Version 3.1 December 2024
% See section 11 of the User Manual for version history
%
%%%%%%%%%%%%%%%%%%%%%%%%%%%%%%%%%%%%%%%%%%%%%%%%%%%%%%%%%%%%%%%%%%%%%%
%%                                                                 %%
%% Please do not use \input{...} to include other tex files.       %%
%% Submit your LaTeX manuscript as one .tex document.              %%
%%                                                                 %%
%% All additional figures and files should be attached             %%
%% separately and not embedded in the \TeX\ document itself.       %%
%%                                                                 %%
%%%%%%%%%%%%%%%%%%%%%%%%%%%%%%%%%%%%%%%%%%%%%%%%%%%%%%%%%%%%%%%%%%%%%

%%\documentclass[referee,sn-basic]{sn-jnl}% referee option is meant for double line spacing

%%=======================================================%%
%% to print line numbers in the margin use lineno option %%
%%=======================================================%%

%%\documentclass[lineno,pdflatex,sn-basic]{sn-jnl}% Basic Springer Nature Reference Style/Chemistry Reference Style

%%=========================================================================================%%
%% the documentclass is set to pdflatex as default. You can delete it if not appropriate.  %%
%%=========================================================================================%%

%%\documentclass[sn-basic]{sn-jnl}% Basic Springer Nature Reference Style/Chemistry Reference Style

%%Note: the following reference styles support Namedate and Numbered referencing. By default the style follows the most common style. To switch between the options you can add or remove “Numbered” in the optional parenthesis. 
%%The option is available for: sn-basic.bst, sn-chicago.bst%  
 
%%\documentclass[pdflatex,sn-nature]{sn-jnl}% Style for submissions to Nature Portfolio journals
\documentclass[pdflatex,sn-basic]{sn-jnl}% Basic Springer Nature Reference Style/Chemistry Reference Style
%%\documentclass[pdflatex,sn-mathphys-num]{sn-jnl}% Math and Physical Sciences Numbered Reference Style
%%\documentclass[pdflatex,sn-mathphys-ay]{sn-jnl}% Math and Physical Sciences Author Year Reference Style
%%\documentclass[pdflatex,sn-aps]{sn-jnl}% American Physical Society (APS) Reference Style
%%\documentclass[pdflatex,sn-vancouver-num]{sn-jnl}% Vancouver Numbered Reference Style
%%\documentclass[pdflatex,sn-vancouver-ay]{sn-jnl}% Vancouver Author Year Reference Style
%%\documentclass[pdflatex,sn-apa]{sn-jnl}% APA Reference Style
%%\documentclass[pdflatex,sn-chicago]{sn-jnl}% Chicago-based Humanities Reference Style

%%%% Standard Packages
%%<additional latex packages if required can be included here>

\usepackage{graphicx}%
\usepackage{multirow}%
\usepackage{amsmath,amssymb,amsfonts}%
\usepackage{amsthm}%
\usepackage{mathrsfs}%
\usepackage[title]{appendix}%
\usepackage{xcolor}%
\usepackage{textcomp}%
\usepackage{manyfoot}%
\usepackage{booktabs}%
\usepackage{algorithm}%
\usepackage{algorithmicx}%
\usepackage{algpseudocode}%
\usepackage{listings}%
\usepackage{booktabs,rotating}
%%%%

%%%%%=============================================================================%%%%
%%%%  Remarks: This template is provided to aid authors with the preparation
%%%%  of original research articles intended for submission to journals published 
%%%%  by Springer Nature. The guidance has been prepared in partnership with 
%%%%  production teams to conform to Springer Nature technical requirements. 
%%%%  Editorial and presentation requirements differ among journal portfolios and 
%%%%  research disciplines. You may find sections in this template are irrelevant 
%%%%  to your work and are empowered to omit any such section if allowed by the 
%%%%  journal you intend to submit to. The submission guidelines and policies 
%%%%  of the journal take precedence. A detailed User Manual is available in the 
%%%%  template package for technical guidance.
%%%%%=============================================================================%%%%

%% as per the requirement new theorem styles can be included as shown below
\theoremstyle{thmstyleone}%
\newtheorem{theorem}{Theorem}%  meant for continuous numbers
%%\newtheorem{theorem}{Theorem}[section]% meant for sectionwise numbers
%% optional argument [theorem] produces theorem numbering sequence instead of independent numbers for Proposition
\newtheorem{proposition}[theorem]{Proposition}% 
%%\newtheorem{proposition}{Proposition}% to get separate numbers for theorem and proposition etc.

\theoremstyle{thmstyletwo}%
\newtheorem{example}{Example}%
\newtheorem{remark}{Remark}%

\theoremstyle{thmstylethree}%
\newtheorem{definition}{Definition}%

\raggedbottom
%%\unnumbered% uncomment this for unnumbered level heads

\begin{document}

\title[Article Title]{Assessing the role of tree mortality in shaping ecosystem functional properties}



%%=============================================================%%
%% GivenName	-> \fnm{Joergen W.}
%% Particle	-> \spfx{van der} -> surname prefix
%% FamilyName	-> \sur{Ploeg}
%% Suffix	-> \sfx{IV}
%% \author*[1,2]{\fnm{Joergen W.} \spfx{van der} \sur{Ploeg} 
%%  \sfx{IV}}\email{iauthor@gmail.com}
%%=============================================================%%

\author*[1,2]{\fnm{Negin} \sur{Katal}}\email{negin.katal@geosense.uni-freiburg.de}

\author[2,3]{\fnm{Second} \sur{Author}}\email{iiauthor@gmail.com}
% \equalcont{These authors contributed equally to this work.}

\author[1,2]{\fnm{Teja} \sur{Kattenborn}}\email{teja.kattenborn@geosense.uni-freiburg.de}
% \equalcont{These authors contributed equally to this work.}

\affil*[1]{\orgdiv{Department}, \orgname{Organization}, \orgaddress{\street{Street}, \city{City}, \postcode{100190}, \state{State}, \country{Country}}}

\affil[2]{\orgdiv{Department}, \orgname{Organization}, \orgaddress{\street{Street}, \city{City}, \postcode{10587}, \state{State}, \country{Country}}}

\affil[3]{\orgdiv{Department}, \orgname{Organization}, \orgaddress{\street{Street}, \city{City}, \postcode{610101}, \state{State}, \country{Country}}}

%%==================================%%
%% Sample for unstructured abstract %%
%%==================================%%

\abstract{EGU Abstract:
Tree mortality is increasing globally due to climate extremes, disturbances, and biotic stressors, with profound consequences for ecosystem carbon cycling. However, the extent to which spatial patterns and temporal dynamics of tree mortality influence the functioning of ecosystems as a whole remains poorly quantified. 
In this study, we investigate how tree mortality impacts  key ecosystem functional properties, focusing on light-saturated gross primary productivity (GPPsat) and maximum net ecosystem production (NEPmax). Ecosystem functional properties were derived from long-term, half-hourly eddy covariance measurements across a range of forest ecosystems. Spatially explicit information on forest cover and tree mortality was obtained from satellite-based predictions of the deadtrees.earth initiative, which integrates drone-based observations with Earth observation data to produce multitemporal mortality and forest cover estimates at global  scale.
We assessed whether including tree mortality information in addition to environmental drives improves the explanation and prediction of global and site-level variability in the flux-derived ecosystem functional properties; GPPsat and NEPmax. Model performance and variable importance patterns were compared between scenarios with and without forest cover and mortality dynamics to quantify the added explanatory power.
This study aims to provide the first systematic assessment of how spatially explicit tree mortality information contributes to ecosystem functional properties derived from eddy covariance data, and to evaluate whether integrating tree mortality observations can improve our understanding of the controls of ecosystem productivity and carbon balance under ongoing climate change.
}


\keywords{ecosystem functioning, tree mortality, eddy covariance, deadtrees.earth}

%%\pacs[JEL Classification]{D8, H51}

%%\pacs[MSC Classification]{35A01, 65L10, 65L12, 65L20, 65L70}

\maketitle

\section{Introduction}\label{sec1}
Terrestrial ecosystems play a fundamental role in regulating global carbon, water, and energy balances through complex interactions between vegetation structure, physiological processes, and environmental drivers (\cite{bonan2008forests, fisher2018vegetation}). Forest ecosystems, in particular, represent critical components of the global climate system, storing approximately 45\% of terrestrial carbon and influencing regional precipitation patterns through evapotranspiration processes (\cite{pan2011large, ellison2017trees}). The functional properties of these ecosystems, including carbon sequestration rates, water use efficiency, and energy partitioning, are intrinsically linked to vegetation characteristics such as leaf area index, canopy structure, and plant functional traits (\cite{wright2004worldwide, diazGlobalSpectrumPlant2016}). However, increasing frequencies of drought, heat stress, and biotic disturbances under changing climatic conditions have led to widespread tree mortality events globally, fundamentally altering ecosystem structure and potentially disrupting these critical ecosystem services (\cite{allen2010global, anderegg2013consequences}).

Tree mortality represents a key ecological process that can dramatically reshape ecosystem functional properties through multiple pathways. Following mortality events, immediate changes occur in canopy structure and leaf area, directly affecting light interception, evapotranspiration rates, and carbon uptake capacity (\cite{hicke2012effects, edburg2012cascading}). These structural changes cascade through the ecosystem, influencing soil microclimate, decomposition rates, and nutrient cycling patterns (\cite{breshears2005regional, adams2009temperature}). Furthermore, the replacement of dead trees by understory vegetation or different plant functional types can lead to long-term shifts in ecosystem functioning, as different vegetation types exhibit distinct physiological characteristics and resource use strategies (\cite{hartmann2018research, hammond2022global}). The magnitude and direction of these changes likely depend on the dominant vegetation type affected, as coniferous and deciduous forests, grasslands, and shrublands exhibit fundamentally different functional trait profiles and responses to disturbance (\cite{franklin1987tree, flower2015responses}).

Despite the recognized importance of tree mortality for ecosystem functioning, significant knowledge gaps persist regarding the quantitative relationships between mortality events and subsequent changes in carbon, water, and energy balances across different vegetation types. Traditional ecological studies have been limited by the spatial and temporal scales of field observations, making it challenging to capture the full spectrum of mortality impacts across diverse ecosystems (\cite{mcdowell2008mechanisms, larson2012tree }). Moreover, the lack of comprehensive datasets linking tree mortality events to plant functional traits has hindered our ability to predict ecosystem responses and develop mechanistic understanding of post-mortality functional changes (\cite{breshears2005regional}). This data limitation has been particularly pronounced for understanding how different plant functional trait combinations influence the magnitude and trajectory of ecosystem recovery following mortality events.

Recent advances in remote sensing technology, machine learning algorithms, and crowd-sourced data collection have opened unprecedented opportunities to address these knowledge gaps at landscape to global scales. The emergence of citizen science platforms and collaborative monitoring networks has dramatically increased the availability of field observations, enabling researchers to document mortality events across diverse ecosystems with greater spatial and temporal resolution than previously possible. In this context, the deadtrees.earth platform represents a significant advancement in mortality monitoring capabilities, providing a standardized framework for collecting and sharing tree mortality data across multiple vegetation types and geographical regions. This platform integrates crowd-sourced field observations with remote sensing data, creating comprehensive datasets that enable large-scale analyses of mortality patterns and their ecosystem consequences.

Building upon these technological advances, recent modeling efforts have begun to incorporate plant functional trait data to better predict ecosystem responses to environmental changes. Notably, Lusk et al. (\cite{lusk2026crowdsourced}) demonstrated how integrating plant functional trait databases with mortality observations can improve predictions of post-disturbance ecosystem functioning, revealing that trait diversity and functional redundancy play critical roles in determining ecosystem resilience following tree mortality events. Their work highlighted the importance of considering plant functional trait composition when assessing vulnerability to mortality and predicting subsequent ecosystem changes, providing a mechanistic framework for understanding functional consequences across different vegetation types.

Given these recent advances in data availability and modeling approaches, there is now an unprecedented opportunity to comprehensively assess the impacts of tree mortality on ecosystem functional properties across diverse vegetation types. The central questions driving this research are: What are the quantitative impacts of tree mortality on ecosystem carbon, water, and energy balances? Do these impacts vary systematically among different vegetation types based on their functional trait composition? And can we identify vegetation types that exhibit greater sensitivity or resilience?
 

\section{Results}\label{sec2}

\subsection{Overview of expected patterns}

We expect ecosystem functional properties (EFPs) to vary substantially across sites and years, reflecting differences in climate, vegetation characteristics, and disturbance intensity. In particular, we expect light-saturated gross primary productivity (GPP$_{sat}$) and maximum net ecosystem production (NEP$_{max}$) to decline with increasing tree mortality and canopy disturbance, although the magnitude of this response will likely vary among sites and forest types.

We further expect that models including forest cover and mortality-related variables will outperform models based only on meteorological and plant trait predictors. This improvement is anticipated to be strongest at sites with marked mortality events or strong temporal variation in canopy condition.

\subsection{Spatial and temporal variation in ecosystem functional properties}

Across the global network of forest eddy covariance sites, annual values of GPP$_{sat}$, NEP$_{max}$, ET$_{max}$, uWUE, and WUE are expected to show pronounced spatial variability. Part of this variability will likely be associated with broad climatic gradients, while another part may reflect structural ecosystem differences captured by community-weighted mean plant traits and forest cover.

We also expect interannual variation within sites, particularly in years affected by disturbance, drought, or canopy decline. Sites experiencing stronger forest cover loss or deadwood accumulation are expected to exhibit reduced productivity and carbon uptake capacity relative to less disturbed years.

\begin{figure}[ht]
\centering
\includegraphics[width=\linewidth]{fig1_site_map_EFP_distributions}
\caption{\textbf{Study sites and ecosystem functional property distributions.}
(\textbf{a}) Global distribution of 184 eddy covariance sites coloured by IGBP land-cover class.
(\textbf{b--e}) Distributions of annual ecosystem functional properties across IGBP classes:
gross primary productivity at light saturation (GPP$_{sat}$),
maximum net ecosystem production (NEP$_{max}$),
maximum evapotranspiration (ET$_{max}$),
underlying water-use efficiency (uWUE), and water-use efficiency (WUE).
Horizontal violins show the full distribution; inner boxplots indicate the interquartile range and median.}
\label{fig:efp_overview}
\end{figure}

\subsection{Patterns in forest cover, mortality, and deadwood metrics}

Forest cover and mortality-related variables are expected to show strong heterogeneity among sites. Some sites will likely show relatively stable canopy conditions, whereas others may exhibit clear year-to-year decreases in forest cover or increases in deadwood fraction within the 500~m tower buffer.

Among the disturbance indicators, we expect the most informative variables to be those capturing both the extent and intensity of canopy loss and deadwood accumulation, such as \texttt{loss\_area\_frac\_500m}, \texttt{loss\_sum\_pp\_500m}, and \texttt{deadwood\_increase\_sum\_pp\_500m}. These variables are expected to provide complementary information about recent mortality dynamics around the flux towers.

\begin{figure}[ht]
\centering
\fbox{\rule{0pt}{2.2in}\rule{0.9\linewidth}{0pt}}
\caption{Placeholder for a figure showing examples of forest cover, loss, and deadwood metrics around selected eddy covariance sites. This figure may later include maps, buffer summaries, or time series of disturbance indicators.}
\label{fig:mortality_metrics}
\end{figure}

\subsection{Expected model performance across predictor sets}

We expect Random Forest models that include only meteorological variables (M1) to explain part of the variation in EFPs, especially for productivity-related variables that are strongly linked to climatic conditions. Adding plant traits (M2) is expected to improve performance by accounting for structural and physiological differences among ecosystems.

The strongest model performance is expected for models that additionally include forest cover and disturbance variables (M3). If tree mortality provides meaningful explanatory information beyond climate and trait variation, then M3 should show higher R$^2$ and lower RMSE than M1 and M2. This improvement is expected to be particularly visible for GPP$_{sat}$ and NEP$_{max}$, which are closely tied to canopy condition and carbon uptake.





%%%%%%%%%%%%%%%%%%%%%%%%%%
\begin{table}[htbp]
\centering
\caption{Leave-one-site-out random forest performance for five ecosystem
functional properties (EFPs), evaluated on v3 harmonized benchmark (1,164 site-years,
164 sites) with complete 24-month predictor data. Each cell shows $R^2$ (RMSE in
native units). \textbf{Bold} = best $R^2$ per EFP within each panel.
Panels A--B: anomaly z-score EFP memory; Panels C--D: raw-lag EFP memory.}
\label{tab:RF_LOSO_v3}
\footnotesize
\begin{tabular}{llcccccc}
\toprule
Model & Predictors & $n$ & GPP$_\mathrm{sat}$ & NEP$_\mathrm{max}$ & ET$_\mathrm{max}$ & uWUE & WUE \\
  &  &  & ($\mu$mol m$^{-2}$ s$^{-1}$) & ($\mu$mol m$^{-2}$ s$^{-1}$) & (mm d$^{-1}$) & (g C mm$^{-1}$) & (g C mm$^{-1}$) \\
\midrule
\multicolumn{8}{l}{\textit{Panel A: 12m window --- Anomaly z-score EFP memory --- 638 site-years, 164 sites}} \\
\midrule
M01 & C & \multirow{8}{*}{638} & 0.367 (8.05) & 0.376 (7.47) & 0.318 (0.064) & 0.162 (1.22) & 0.247 (1.17) \\
M02 & C + D &  & 0.542 (6.86) & 0.518 (6.57) & 0.346 (0.062) & \textbf{0.212 (1.17)} & 0.298 (1.13) \\
M03 & C + T &  & 0.505 (7.11) & 0.500 (6.67) & \textbf{0.476 (0.056)} & 0.178 (1.20) & 0.241 (1.18) \\
M04 & C + T + D &  & \textbf{0.573 (6.63)} & \textbf{0.555 (6.31)} & 0.461 (0.057) & 0.197 (1.18) & 0.263 (1.16) \\
M05 & C + M &  & 0.362 (8.08) & 0.363 (7.54) & 0.316 (0.064) & 0.161 (1.22) & 0.235 (1.18) \\
M06 & C + D + M &  & 0.544 (6.85) & 0.510 (6.62) & 0.343 (0.062) & 0.205 (1.18) & \textbf{0.306 (1.12)} \\
M07 & C + T + M &  & 0.506 (7.10) & 0.503 (6.65) & 0.462 (0.057) & 0.177 (1.20) & 0.248 (1.17) \\
M08 & C + T + D + M &  & 0.569 (6.66) & 0.548 (6.35) & 0.458 (0.057) & 0.187 (1.19) & 0.268 (1.15) \\
\midrule
\multicolumn{8}{l}{\textit{Panel B: 24m window --- Anomaly z-score EFP memory --- 638 site-years, 164 sites}} \\
\midrule
M01 & C & \multirow{8}{*}{638} & 0.374 (8.00) & 0.370 (7.50) & 0.325 (0.063) & 0.169 (1.21) & 0.261 (1.16) \\
M02 & C + D &  & 0.537 (6.90) & 0.508 (6.64) & 0.332 (0.063) & \textbf{0.199 (1.18)} & 0.280 (1.14) \\
M03 & C + T &  & 0.509 (7.09) & 0.505 (6.63) & \textbf{0.473 (0.056)} & 0.181 (1.20) & 0.256 (1.17) \\
M04 & C + T + D &  & \textbf{0.572 (6.65)} & 0.546 (6.37) & 0.457 (0.057) & 0.198 (1.18) & 0.262 (1.16) \\
M05 & C + M &  & 0.372 (8.00) & 0.371 (7.48) & 0.309 (0.064) & 0.166 (1.21) & 0.242 (1.17) \\
M06 & C + D + M &  & 0.538 (6.90) & 0.505 (6.66) & 0.324 (0.063) & 0.198 (1.18) & \textbf{0.282 (1.14)} \\
M07 & C + T + M &  & 0.502 (7.13) & 0.508 (6.61) & 0.467 (0.056) & 0.172 (1.20) & 0.258 (1.16) \\
M08 & C + T + D + M &  & 0.571 (6.65) & \textbf{0.549 (6.35)} & 0.449 (0.057) & 0.186 (1.19) & 0.247 (1.17) \\
\midrule
\multicolumn{8}{l}{\textit{Panel C: 12m window --- Raw-lag EFP memory --- 631 site-years, 164 sites}} \\
\midrule
M01 & C & \multirow{8}{*}{631} & 0.359 (8.10) & 0.363 (7.56) & 0.309 (0.064) & 0.162 (1.22) & 0.236 (1.18) \\
M02 & C + D &  & 0.539 (6.88) & 0.513 (6.61) & 0.351 (0.062) & 0.207 (1.18) & 0.294 (1.13) \\
M03 & C + T &  & 0.505 (7.11) & 0.508 (6.62) & 0.471 (0.056) & 0.168 (1.21) & 0.244 (1.18) \\
M04 & C + T + D &  & 0.568 (6.66) & 0.546 (6.37) & 0.451 (0.057) & 0.182 (1.19) & 0.251 (1.17) \\
M05 & C + M &  & \textbf{0.821 (4.29)} & \textbf{0.842 (3.78)} & \textbf{0.704 (0.042)} & \textbf{0.604 (0.837)} & \textbf{0.601 (0.859)} \\
M06 & C + D + M &  & 0.792 (4.76) & 0.795 (4.43) & 0.639 (0.048) & 0.510 (0.939) & 0.528 (0.943) \\
M07 & C + T + M &  & 0.802 (4.59) & 0.816 (4.16) & 0.656 (0.046) & 0.561 (0.894) & 0.572 (0.903) \\
M08 & C + T + D + M &  & 0.777 (4.92) & 0.779 (4.58) & 0.621 (0.049) & 0.477 (0.970) & 0.501 (0.972) \\
\midrule
\multicolumn{8}{l}{\textit{Panel D: 24m window --- Raw-lag EFP memory --- 631 site-years, 164 sites}} \\
\midrule
M01 & C & \multirow{8}{*}{631} & 0.375 (7.99) & 0.368 (7.52) & 0.316 (0.064) & 0.164 (1.22) & 0.251 (1.17) \\
M02 & C + D &  & 0.538 (6.89) & 0.499 (6.70) & 0.338 (0.063) & 0.194 (1.19) & 0.280 (1.15) \\
M03 & C + T &  & 0.504 (7.12) & 0.508 (6.62) & 0.467 (0.056) & 0.164 (1.21) & 0.245 (1.17) \\
M04 & C + T + D &  & 0.569 (6.66) & 0.550 (6.34) & 0.458 (0.057) & 0.189 (1.19) & 0.253 (1.17) \\
M05 & C + M &  & \textbf{0.833 (4.14)} & \textbf{0.851 (3.67)} & \textbf{0.733 (0.040)} & \textbf{0.604 (0.835)} & \textbf{0.598 (0.860)} \\
M06 & C + D + M &  & 0.818 (4.43) & 0.814 (4.20) & 0.684 (0.045) & 0.517 (0.930) & 0.538 (0.930) \\
M07 & C + T + M &  & 0.827 (4.27) & 0.836 (3.89) & 0.713 (0.042) & 0.571 (0.875) & 0.580 (0.888) \\
M08 & C + T + D + M &  & 0.810 (4.52) & 0.810 (4.24) & 0.665 (0.046) & 0.507 (0.941) & 0.527 (0.943) \\
\midrule
\bottomrule
\end{tabular}
\vspace{2pt}
\begin{minipage}{\linewidth}
\footnotesize
\textit{Note:} C = Climate (monthly TA, VPD, SW$_\mathrm{IN}$ mean/p05/p95;
P mean, sum, $P_\mathrm{p05}$, $P_\mathrm{p95}$); T = Plant traits (hydraulic, leaf, root);
D = Disturbance (mortality intensity, deadwood, forest loss at 100--500\,m);
M = EFP memory (anomaly z-score in Panels A--B; raw lag-1/lag-2 values in Panels C--D).
RMSE in native units in parentheses. Data quality v2: negative uWUE/WUE set to NA;
CN-Sb1/CN-Sb2 extreme outliers removed; US-HB4 anomalous precipitation months set to NA.
\end{minipage}
\end{table}


%%%%%%%%%%%%%%%%%%%%%%%%%


\subsection{Expected contribution of tree mortality variables}

We expect the inclusion of mortality-related predictors to increase model explanatory power, although the magnitude of improvement will likely differ across response variables. The strongest contribution is expected for GPP$_{sat}$ and NEP$_{max}$, because tree mortality directly alters canopy structure, photosynthetic capacity, and stand-level carbon balance.

Variable importance analyses are expected to show that meteorological predictors remain dominant overall, but that mortality-related variables become particularly influential at disturbed sites or in years following canopy loss. This would support the hypothesis that tree mortality contributes meaningful additional information for explaining ecosystem functioning beyond climate and plant traits alone.

\begin{figure}[ht]
\centering
\fbox{\rule{0pt}{2.2in}\rule{0.9\linewidth}{0pt}}
\caption{Placeholder for variable importance plots comparing meteorological, trait, and mortality predictors across ecosystem functional properties.}
\label{fig:varimp}
\end{figure}

\subsection{Expected site-level heterogeneity}

The effect of tree mortality is not expected to be uniform across all sites. Instead, we anticipate strong site-level heterogeneity depending on background climate, vegetation type, disturbance history, and resilience capacity. Some sites may show clear reductions in GPP$_{sat}$ or NEP$_{max}$ following canopy loss, whereas other sites may exhibit weaker or delayed responses.

This heterogeneity is expected to highlight that the role of mortality depends not only on the amount of canopy loss itself, but also on the broader environmental and ecological context in which mortality occurs.

\begin{figure}[ht]
\centering
\fbox{\rule{0pt}{2.2in}\rule{0.9\linewidth}{0pt}}
\caption{Placeholder for site-level model evaluation results, for example leave-one-site-out performance or site-specific response patterns to tree mortality.}
\label{fig:site_level}
\end{figure}

%%%%%%%%


\subsection{Interim interpretation}

Taken together, the expected results support the working hypothesis that tree mortality is an important but spatially variable control on ecosystem functioning. While meteorological conditions and plant functional traits are expected to explain a substantial share of variation in EFPs, forest cover and mortality dynamics are likely to provide additional explanatory power, especially for productivity- and carbon-balance-related properties.

These anticipated findings would suggest that integrating remotely sensed mortality information into flux-based analyses can improve our understanding of how forest disturbance shapes ecosystem functioning under ongoing environmental change.


%%=============================================%%
%% For presentation purpose, we have included  %%
%% \bigskip command. Please ignore this.       %%
%%=============================================%%
% \bigskip
% \begin{verbatim}
% \begin{table}[<placement-specifier>]
% \caption{<table-caption>}\label{<table-label>}%
% \begin{tabular}{@{}llll@{}}
% \toprule
% Column 1 & Column 2 & Column 3 & Column 4\\
% \midrule
% row 1 & data 1 & data 2	 & data 3 \\
% row 2 & data 4 & data 5\footnotemark[1] & data 6 \\
% row 3 & data 7 & data 8	 & data 9\footnotemark[2]\\
% \botrule
% \end{tabular}
% \footnotetext{Source: This is an example of table footnote. 
% This is an example of table footnote.}
% \footnotetext[1]{Example for a first table footnote.
% This is an example of table footnote.}
% \footnotetext[2]{Example for a second table footnote. 
% This is an example of table footnote.}
% \end{table}
% \end{verbatim}
% \bigskip
%%=============================================%%
%% For presentation purpose, we have included  %%
%% \bigskip command. Please ignore this.       %%
%%=============================================%%



\section{Methods}\label{sec11}

\subsection*{Study sites and data sources}

We assembled a multi-network eddy covariance dataset spanning 166 sites and 644 site-years across the period 2017–2025, drawing from FLUXNET, ICOS, AmeriFlux, NEON, and OzFlux. Sites were included if they met two criteria: (i) at least four consecutive years of valid half-hourly flux observations within the study window, and (ii) availability of co-located forest cover and disturbance data from the deadtrees.earth platform \cite{mosig2026deadtrees}. The final dataset covers a range of IGBP land-cover classes including evergreen needleleaf forest (ENF; $n=58$ sites), deciduous broadleaf forest (DBF; $n=31$), mixed forest (MF; $n=7$), and open ecosystems including shrublands, savannas, and wetlands (OSH, WSA, SAV, CSH, WET; $n=70$).

\subsection*{Ecosystem functional properties}

Annual ecosystem functional properties (EFPs) were derived from half-hourly eddy covariance flux observations following the methodology of \cite{musavi2015imprint} and \cite{migliavacca2021three}. For each site-year, we computed:
\begin{itemize}
    \item Light-saturated gross primary productivity (\textbf{GPP$_\mathbf{sat}$}; $\mu$mol m$^{-2}$ s$^{-1}$): the upper quantile of the light-response curve fitted to half-hourly GPP observations,
    \item Maximum net ecosystem production (\textbf{NEP$_\mathbf{max}$}; $\mu$mol m$^{-2}$ s$^{-1}$): the upper-envelope carbon uptake capacity,
    \item Maximum evapotranspiration (\textbf{ET$_\mathbf{max}$}; mm d$^{-1}$): the upper-envelope water flux,
    \item Underlying water-use efficiency (\textbf{uWUE}; g C mm$^{-1}$): defined as $\text{GPP} \cdot \sqrt{\text{VPD}} / \text{ET}$, which accounts for atmospheric demand \cite{migliavacca_three_2021},
    \item Water-use efficiency (\textbf{WUE}; g C mm$^{-1}$): defined as GPP / ET.
\end{itemize}

EFP estimation was anchored to the center of the growing season (CGS), defined as the date of maximum daily GPP after smoothing. This approach avoids phenological confounding and ensures predictors reflect antecedent conditions relative to peak carbon uptake.

\subsection*{Meteorological predictors}

Meteorological drivers were computed from on-site eddy covariance observations. For four variables --- vapor pressure deficit (VPD), air temperature (T$_\text{air}$), precipitation (P), and incoming shortwave radiation (SW$_\text{in}$) --- we derived monthly aggregates (mean, 5th and 95th percentiles for VPD and T$_\text{air}$; mean and cumulative sum for P; mean for SW$_\text{in}$) for each of the 24 calendar months preceding the CGS peak. This produced a set of time-lagged meteorological predictors capturing both seasonal and inter-annual climatic memory, with two window lengths evaluated: the 12 months immediately preceding the CGS (12m), and the full 24-month window (24m). The 24-month configuration was used as the primary benchmark.

\subsection*{Plant functional trait data}

Site-level community-weighted mean (CWM) plant functional traits were obtained from the planttraits.earth database \cite{lusk2026crowdsourced}. Fourteen trait variables were included, covering leaf morphology (specific leaf area, leaf area, leaf thickness, leaf dry mass), leaf chemistry (leaf N and P content by mass and area, leaf C, leaf C:N ratio, leaf $\delta^{15}$N), and rooting depth. These traits represent time-invariant structural and physiological characteristics and were used to capture baseline differences in ecosystem functioning independent of recent disturbance.

\subsection*{Forest cover and disturbance metrics}

Spatially explicit annual forest cover and deadwood fraction maps were derived from the Sentinel-2--based deadtrees.earth platform \cite{mosig2026deadtrees} using version 2-2 of the disturbance product. For each flux tower, circular buffers of five radii (100, 200, 300, 400, and 500 m) were delineated in the local UTM projection. Within each buffer, the following quantities were computed for each year:

\begin{itemize}
    \item \textit{Forest cover fraction} (\texttt{forest\_mean\_pct}): mean canopy cover across all buffer pixels (\%),
    \item \textit{Deadwood fraction} (\texttt{deadwood\_mean\_pct}): mean fraction of pixels classified as standing deadwood (\%),
    \item \textit{Forest loss area} (\texttt{loss\_area\_frac}): fraction of pixels experiencing a year-on-year decrease in forest cover exceeding 20 percentage points,
    \item \textit{Forest loss magnitude} (\texttt{loss\_mean\_pp}, \texttt{loss\_sum\_pp}): mean and total magnitude of forest cover loss across affected pixels,
    \item \textit{Deadwood increase area} (\texttt{deadwood\_increase\_area\_frac}): fraction of pixels with a year-on-year increase in deadwood exceeding 20 percentage points,
    \item \textit{Deadwood increase magnitude} (\texttt{deadwood\_increase\_mean\_pp}, \texttt{deadwood\_increase\_sum\_pp}): mean and total magnitude of deadwood increase.
\end{itemize}

For each of these seven per-buffer metrics, we also included one- and two-year lagged values, yielding a total of 130 disturbance predictor variables. The 500 m buffer was used for all main analyses, as it best represents the flux tower footprint for the majority of sites. The \texttt{deadwood\_mean\_pct\_500m} variable --- representing the cumulative standing deadwood fraction --- emerged as the dominant disturbance signal in subsequent SHAP analyses and is used as the primary site-level disturbance descriptor throughout.

\subsection*{Modeling framework}

We trained Random Forest models \cite{breiman2001random} using the \texttt{ranger} implementation in R \cite{wright2017ranger} to predict annual EFPs from combinations of meteorological, trait, and disturbance predictors. Each forest consisted of 500 trees. To disentangle the contributions of different predictor groups, we defined eight model configurations (M01--M08) by systematically including or excluding three predictor blocks:

\begin{table}[h!]
\centering
\small
\begin{tabular}{lcccc}
\hline
Model & Climate (C) & Traits (T) & Disturbance (D) & EFP memory (M) \\
\hline
M01 & \checkmark & & & \\
M02 & \checkmark & & \checkmark & \\
M03 & \checkmark & \checkmark & & \\
M04 & \checkmark & \checkmark & \checkmark & \\
M05 & \checkmark & & & \checkmark \\
M06 & \checkmark & & \checkmark & \checkmark \\
M07 & \checkmark & \checkmark & & \checkmark \\
M08 & \checkmark & \checkmark & \checkmark & \checkmark \\
\hline
\end{tabular}
\caption{Predictor block combinations for the eight model configurations. C = monthly meteorological aggregates; T = community-weighted mean plant traits; D = remote sensing--based disturbance and deadwood metrics; M = lagged EFP (ecosystem memory).}
\label{tab:models}
\end{table}

EFP memory (M) was encoded in two ways: (i) as the z-score anomaly of the lagged EFP relative to the site-level mean (anomaly encoding, primary), and (ii) as the raw lagged EFP value (raw-lag encoding, presented as a sensitivity comparison). Models M01--M04 therefore operate without memory, while M05--M08 additionally condition on the previous year's ecosystem state.

All 32 configurations (8 model structures $\times$ 2 memory types $\times$ 2 meteorological windows) were trained and evaluated independently for each of the five EFPs.

\subsection*{Model validation}

Models were validated using leave-one-site-out cross-validation (LOSO-CV), in which all observations from one site are withheld as a test set while the remaining 165 sites form the training set. This procedure was repeated for each of the 166 sites, yielding out-of-sample predictions for all 644 site-years. LOSO-CV preserves spatial independence between training and test sets, providing a conservative estimate of model transferability to unseen sites \cite{roberts2017cross}. Performance was quantified using the root mean square error (RMSE) and coefficient of determination (R$^2$).

\subsection*{Quantifying the effect of disturbance}

The contribution of disturbance predictors was assessed by comparing paired model configurations that differ only in the inclusion of the disturbance block (D): M01 vs.\ M02, M03 vs.\ M04, M05 vs.\ M06, and M07 vs.\ M08. For each site and model pair, we computed the per-site $\Delta$RMSE as:

\begin{equation}
\Delta\text{RMSE} = \text{RMSE}_{\text{with D}} - \text{RMSE}_{\text{without D}}
\end{equation}

Negative values indicate that adding disturbance predictors improved prediction accuracy; positive values indicate degraded performance. The primary comparison was M03 vs.\ M04 (C+T vs.\ C+T+D), which represents the hardest test: disturbance must improve upon a model already informed by both climate and plant traits.

To characterize the spatial heterogeneity of the disturbance effect, we stratified sites into three deadwood tiers based on their mean \texttt{deadwood\_mean\_pct\_500m}: Low ($<$5\%, $n=58$ sites), Medium (5--12\%, $n=83$ sites), and High ($>$12\%, $n=43$ sites). Additional stratification by IGBP land-cover class was used to assess biome-level differences in the disturbance response.

\subsection*{Variable importance and SHAP analysis}

To interpret which disturbance variables drive model predictions, we computed SHapley Additive exPlanations (SHAP) values \cite{lundberg2017unified} using the \texttt{treeshap} package \cite{kozmarski2021treeshap} for models M04 and M08 across both memory encodings and lag windows. SHAP values quantify the marginal contribution of each predictor to each prediction, averaged over all possible predictor orderings. Predictors were grouped into four blocks (Climate, Traits, Disturbance, Memory) and grouped SHAP magnitudes were used to assess the relative importance of each block.




% \begin{enumerate}[1.]
% \item Approval: a statement which confirms that all experimental protocols were approved by a named institutional and/or licensing committee. Please identify the approving body in the methods section

% \item Accordance: a statement explicitly saying that the methods were carried out in accordance with the relevant guidelines and regulations

% \item Informed consent (for experiments involving humans or human tissue samples): include a statement confirming that informed consent was obtained from all participants and/or their legal guardian/s
% \end{enumerate}


\section{Discussion}\label{sec12}


\section{Conclusion}\label{sec13}



\backmatter

\bmhead{Supplementary information}

If your article has accompanying supplementary file/s please state so here. 

Authors reporting data from electrophoretic gels and blots should supply the full unprocessed scans for key as part of their Supplementary information. This may be requested by the editorial team/s if it is missing.

Please refer to Journal-level guidance for any specific requirements.

\bmhead{Acknowledgements}



\section*{Declarations}



\begin{itemize}
\item Funding
\item Conflict of interest/Competing interests (check journal-specific guidelines for which heading to use)
\item Ethics approval and consent to participate
\item Consent for publication
\item Data availability 
\item Materials availability
\item Code availability 
\item Author contribution
\end{itemize}

\noindent
If any of the sections are not relevant to your manuscript, please include the heading and write `Not applicable' for that section. 

%%===================================================%%
%% For presentation purpose, we have included        %%
%% \bigskip command. Please ignore this.             %%
%%===================================================%%
\bigskip
\begin{flushleft}%
Editorial Policies for:

\bigskip\noindent
Springer journals and proceedings: \url{https://www.springer.com/gp/editorial-policies}

\bigskip\noindent
Nature Portfolio journals: \url{https://www.nature.com/nature-research/editorial-policies}

\bigskip\noindent
\textit{Scientific Reports}: \url{https://www.nature.com/srep/journal-policies/editorial-policies}

\bigskip\noindent
BMC journals: \url{https://www.biomedcentral.com/getpublished/editorial-policies}
\end{flushleft}

\begin{appendices}

\section{Section title of first appendix}\label{secA1}

An appendix contains supplementary information that is not an essential part of the text itself but which may be helpful in providing a more comprehensive understanding of the research problem or it is information that is too cumbersome to be included in the body of the paper.

%%=============================================%%
%% For submissions to Nature Portfolio Journals %%
%% please use the heading ``Extended Data''.   %%
%%=============================================%%

%%=============================================================%%
%% Sample for another appendix section			       %%
%%=============================================================%%

%% \section{Example of another appendix section}\label{secA2}%
%% Appendices may be used for helpful, supporting or essential material that would otherwise 
%% clutter, break up or be distracting to the text. Appendices can consist of sections, figures, 
%% tables and equations etc.

\end{appendices}

%%===========================================================================================%%
%% If you are submitting to one of the Nature Portfolio journals, using the eJP submission   %%
%% system, please include the references within the manuscript file itself. You may do this  %%
%% by copying the reference list from your .bbl file, paste it into the main manuscript .tex %%
%% file, and delete the associated \verb+\bibliography+ commands.                            %%
%%===========================================================================================%%

\bibliography{sn-bibliography}% common bib file
%% if required, the content of .bbl file can be included here once bbl is generated
%%\input sn-article.bbl




%%%%%%%%%%%%%%%%%%%%%%%%%%%%%%%%%%%%%%%%%%
%% Appendix
%%%%%%%%%%%%%%%%%%%%%%%%%%%%%%%%%%%%%%%%%%
\appendix

\section{Supplementary Tables}\label{sec:appendix-tables}

\setcounter{table}{0}
\renewcommand{\thetable}{A\arabic{table}}

%% Table A1: RMSE ratio per biome × EFP
\begin{table}[htbp]
\centering
\caption{RMSE ratio (RMSE with D / RMSE without D) per biome and EFP
(v2 data, anomaly memory, 24-month window). \textbf{Bold} $<1$: D improved;
\textit{italic} $>1$: D worsened. `---' = no data after quality control.}
\label{tab:rmse_ratio_v2}
\begin{tabular}{@{}p{0.48\linewidth}@{\hspace{0.04\linewidth}}p{0.48\linewidth}@{}}
\footnotesize
\setlength{\tabcolsep}{3.5pt}
\renewcommand{\arraystretch}{1.06}
\begin{tabular}{lccccc}
\toprule
\multicolumn{6}{l}{\textit{Panel A: C vs C$+$D}} \\
\midrule
Biome & GPP$_\mathrm{sat}$ & NEP$_\mathrm{max}$ & ET$_\mathrm{max}$ & uWUE & WUE \\
\midrule
ENF & \textbf{0.843} & \textbf{0.950} & \textbf{0.939} & \textbf{0.873} & \textbf{0.815} \\
EBF & \textit{1.217} & \textit{1.133} & \textbf{0.902} & \textit{1.070} & \textit{1.217} \\
DNF & --- & --- & --- & --- & --- \\
DBF & \textbf{0.979} & \textbf{0.953} & \textit{1.121} & \textbf{0.817} & \textbf{0.918} \\
MF & \textbf{0.614} & \textbf{0.559} & \textbf{0.826} & \textbf{0.832} & \textbf{0.731} \\
CSH & \textbf{0.779} & \textbf{0.730} & \textbf{0.917} & \textit{1.156} & \textbf{0.861} \\
OSH & \textbf{0.928} & \textbf{0.859} & \textbf{0.921} & \textit{1.098} & \textit{1.082} \\
WSA & \textbf{0.936} & \textit{1.141} & \textit{1.117} & \textit{1.105} & \textit{1.130} \\
SAV & \textbf{0.854} & \textbf{0.901} & \textit{1.170} & \textit{1.048} & \textbf{0.884} \\
WET & \textbf{0.928} & \textbf{0.942} & \textbf{0.965} & \textbf{0.906} & \textbf{0.974} \\
\midrule
\textit{All} & \textbf{0.901} & \textbf{0.925} & \textbf{0.982} & \textbf{0.958} & \textbf{0.959} \\
\bottomrule
\end{tabular} &
\footnotesize
\setlength{\tabcolsep}{3.5pt}
\renewcommand{\arraystretch}{1.06}
\begin{tabular}{lccccc}
\toprule
\multicolumn{6}{l}{\textit{Panel B: C$+$T vs C$+$T$+$D}} \\
\midrule
Biome & GPP$_\mathrm{sat}$ & NEP$_\mathrm{max}$ & ET$_\mathrm{max}$ & uWUE & WUE \\
\midrule
ENF & \textbf{0.897} & \textbf{0.948} & \textit{1.006} & \textbf{0.911} & \textbf{0.878} \\
EBF & \textit{1.333} & \textit{1.089} & \textbf{0.919} & \textbf{0.949} & \textbf{0.951} \\
DNF & --- & --- & --- & --- & --- \\
DBF & \textbf{0.966} & \textbf{0.973} & \textit{1.056} & \textbf{0.928} & \textbf{0.958} \\
MF & \textbf{0.821} & \textbf{0.801} & \textit{1.002} & \textit{1.050} & \textit{1.216} \\
CSH & \textit{1.265} & \textit{1.283} & \textit{1.012} & \textit{1.314} & \textit{1.055} \\
OSH & \textbf{0.944} & \textbf{0.953} & \textit{1.038} & \textit{1.038} & \textit{1.040} \\
WSA & \textbf{0.989} & \textit{1.033} & \textbf{0.929} & \textit{1.049} & \textit{1.094} \\
SAV & \textit{1.068} & \textit{1.058} & \textit{1.112} & \textit{1.031} & \textit{1.025} \\
WET & \textbf{0.989} & \textbf{0.990} & \textit{1.079} & \textbf{0.969} & \textit{1.013} \\
\midrule
\textit{All} & \textbf{0.956} & \textbf{0.971} & \textit{1.037} & \textbf{0.983} & \textbf{0.987} \\
\bottomrule
\end{tabular} \\[1.4em]
\footnotesize
\setlength{\tabcolsep}{3.5pt}
\renewcommand{\arraystretch}{1.06}
\begin{tabular}{lccccc}
\toprule
\multicolumn{6}{l}{\textit{Panel C: C$+$M vs C$+$D$+$M}} \\
\midrule
Biome & GPP$_\mathrm{sat}$ & NEP$_\mathrm{max}$ & ET$_\mathrm{max}$ & uWUE & WUE \\
\midrule
ENF & \textbf{0.866} & \textbf{0.980} & \textbf{0.938} & \textbf{0.894} & \textbf{0.830} \\
EBF & \textit{1.080} & \textit{1.056} & \textbf{0.861} & \textit{1.022} & \textit{1.149} \\
DNF & --- & --- & --- & --- & --- \\
DBF & \textbf{0.943} & \textbf{0.924} & \textit{1.102} & \textbf{0.856} & \textbf{0.922} \\
MF & \textbf{0.618} & \textbf{0.569} & \textbf{0.882} & \textbf{0.956} & \textbf{0.830} \\
CSH & \textbf{0.791} & \textbf{0.745} & \textbf{0.964} & \textit{1.150} & \textbf{0.894} \\
OSH & \textbf{0.943} & \textbf{0.846} & \textbf{0.959} & \textit{1.109} & \textit{1.067} \\
WSA & \textbf{0.962} & \textit{1.097} & \textit{1.084} & \textit{1.087} & \textit{1.160} \\
SAV & \textbf{0.957} & \textit{1.019} & \textit{1.238} & \textit{1.025} & \textbf{0.839} \\
WET & \textbf{0.919} & \textbf{0.958} & \textbf{0.982} & \textbf{0.902} & \textbf{0.941} \\
\midrule
\textit{All} & \textbf{0.902} & \textbf{0.934} & \textbf{0.994} & \textbf{0.963} & \textbf{0.948} \\
\bottomrule
\end{tabular} &
\footnotesize
\setlength{\tabcolsep}{3.5pt}
\renewcommand{\arraystretch}{1.06}
\begin{tabular}{lccccc}
\toprule
\multicolumn{6}{l}{\textit{Panel D: C$+$T$+$M vs C$+$T$+$D$+$M}} \\
\midrule
Biome & GPP$_\mathrm{sat}$ & NEP$_\mathrm{max}$ & ET$_\mathrm{max}$ & uWUE & WUE \\
\midrule
ENF & \textbf{0.888} & \textbf{0.955} & \textit{1.008} & \textbf{0.914} & \textbf{0.882} \\
EBF & \textit{1.189} & \textit{1.018} & \textbf{0.844} & \textbf{0.973} & \textbf{0.966} \\
DNF & --- & --- & --- & --- & --- \\
DBF & \textbf{0.958} & \textbf{0.959} & \textit{1.044} & \textbf{0.938} & \textbf{0.973} \\
MF & \textbf{0.778} & \textbf{0.823} & \textbf{0.966} & \textit{1.091} & \textit{1.214} \\
CSH & \textit{1.321} & \textit{1.317} & \textit{1.081} & \textit{1.311} & \textit{1.071} \\
OSH & \textbf{0.961} & \textbf{0.935} & \textit{1.064} & \textit{1.049} & \textit{1.055} \\
WSA & \textbf{0.970} & \textit{1.005} & \textbf{0.922} & \textit{1.053} & \textit{1.140} \\
SAV & \textit{1.025} & \textit{1.086} & \textit{1.112} & \textit{1.058} & \textbf{0.997} \\
WET & \textbf{0.983} & \textbf{0.995} & \textit{1.081} & \textbf{0.963} & \textit{1.024} \\
\midrule
\textit{All} & \textbf{0.949} & \textbf{0.973} & \textit{1.038} & \textbf{0.988} & \textbf{0.998} \\
\bottomrule
\end{tabular} \\
\end{tabular}
\end{table}


%% Table A2: LOSO performance — wetland sites excluded
\begin{sidewaystable}[p]
\centering
\caption{Leave-one-site-out cross-validation $R^2$ (RMSE) of random forest
models predicting four ecosystem functional properties (EFPs) on the unified
24-month benchmark. \textbf{Bold}: best $R^2$ per EFP and column group within
each panel. Left group: all sites (644 site-years, 166 sites); right group:
wetland sites excluded (515 site-years, 128 sites). Panels A--B: anomaly
$z$-score EFP memory (M05--M08); Panels C--D: raw-lag EFP memory.
C\,=\,Climate; T\,=\,Traits; D\,=\,Disturbance; M\,=\,EFP memory.
RMSE units: $\mu$mol\,m$^{-2}$\,s$^{-1}$ (GPPsat, NEPmax),
mm\,d$^{-1}$ (ETmax), g\,C\,mm$^{-1}$ (uWUE).}
\label{tab:RF_LOSO_24mbench_noWET}
\scriptsize
\setlength{\tabcolsep}{3pt}
\renewcommand{\arraystretch}{1.0}
\begin{tabular}{@{}llcccccccc@{}}
\toprule
& & \multicolumn{4}{c}{\textbf{All sites} ($n=644$, 166 sites)} &
      \multicolumn{4}{c}{\textbf{Excl.\ wetlands} ($n=515$, 128 sites)} \\
\cmidrule(lr){3-6}\cmidrule(lr){7-10}
Model & Predictors & GPPsat & NEPmax & ETmax & uWUE &
                     GPPsat & NEPmax & ETmax & uWUE \\
\midrule
\multicolumn{10}{l}{\textit{Panel\,A: 12-month window --- anomaly $z$-score memory}} \\
\midrule
M01 & C & 0.338\,(8.3) & 0.328\,(7.9) & 0.311\,(0.064) & 0.102\,(1.67) & 0.412\,(7.8) & 0.380\,(7.6) & 0.297\,(0.062) & 0.067\,(1.69) \\
M02 & C+D & 0.519\,(7.1) & 0.483\,(7.0) & 0.339\,(0.062) & 0.146\,(1.61) & 0.582\,(6.6) & 0.520\,(6.7) & 0.303\,(0.061) & 0.081\,(1.64) \\
M03 & C+T & 0.501\,(7.2) & 0.497\,(6.8) & \textbf{0.459\,(0.057)} & 0.189\,(1.57) & 0.557\,(6.9) & 0.554\,(6.5) & \textbf{0.408\,(0.057)} & 0.160\,(1.58) \\
M04 & C+T+D & 0.555\,(6.8) & 0.530\,(6.6) & 0.450\,(0.057) & \textbf{0.198\,(1.56)} & 0.604\,(6.5) & 0.570\,(6.4) & 0.391\,(0.058) & 0.160\,(1.57) \\
M05 & C+M & 0.327\,(8.4) & 0.320\,(8.0) & 0.307\,(0.064) & 0.081\,(1.69) & 0.406\,(7.9) & 0.369\,(7.7) & 0.299\,(0.062) & 0.055\,(1.70) \\
M06 & C+D+M & 0.520\,(7.1) & 0.489\,(6.9) & 0.334\,(0.063) & 0.142\,(1.62) & 0.583\,(6.6) & 0.530\,(6.7) & 0.294\,(0.062) & 0.081\,(1.64) \\
M07 & C+T+M & 0.504\,(7.2) & 0.499\,(6.8) & 0.458\,(0.057) & 0.179\,(1.58) & 0.559\,(6.9) & 0.552\,(6.5) & 0.408\,(0.057) & \textbf{0.165\,(1.57)} \\
M08 & C+T+D+M & \textbf{0.561\,(6.8)} & \textbf{0.537\,(6.6)} & 0.454\,(0.057) & 0.194\,(1.57) & \textbf{0.609\,(6.4)} & \textbf{0.580\,(6.3)} & 0.399\,(0.057) & 0.154\,(1.57) \\
\midrule
\multicolumn{10}{l}{\textit{Panel\,B: 24-month window --- anomaly $z$-score memory}} \\
\midrule
M01 & C & 0.344\,(8.3) & 0.326\,(8.0) & 0.318\,(0.063) & 0.102\,(1.67) & 0.414\,(7.8) & 0.362\,(7.8) & 0.305\,(0.062) & 0.069\,(1.69) \\
M02 & C+D & 0.518\,(7.1) & 0.472\,(7.0) & 0.325\,(0.063) & 0.118\,(1.64) & 0.584\,(6.6) & 0.510\,(6.8) & 0.281\,(0.062) & 0.067\,(1.66) \\
M03 & C+T & 0.502\,(7.2) & 0.490\,(6.9) & \textbf{0.461\,(0.056)} & 0.193\,(1.57) & 0.566\,(6.8) & 0.541\,(6.6) & \textbf{0.413\,(0.057)} & 0.167\,(1.57) \\
M04 & C+T+D & 0.559\,(6.8) & 0.533\,(6.6) & 0.441\,(0.058) & 0.205\,(1.56) & 0.611\,(6.4) & 0.572\,(6.4) & 0.384\,(0.058) & 0.168\,(1.56) \\
M05 & C+M & 0.344\,(8.3) & 0.323\,(8.0) & 0.304\,(0.064) & 0.092\,(1.68) & 0.422\,(7.8) & 0.365\,(7.7) & 0.297\,(0.062) & 0.066\,(1.68) \\
M06 & C+D+M & 0.517\,(7.1) & 0.477\,(7.0) & 0.319\,(0.063) & 0.121\,(1.64) & 0.581\,(6.7) & 0.516\,(6.8) & 0.273\,(0.063) & 0.067\,(1.65) \\
M07 & C+T+M & 0.504\,(7.2) & 0.496\,(6.8) & 0.457\,(0.057) & \textbf{0.213\,(1.55)} & 0.560\,(6.9) & 0.548\,(6.6) & 0.404\,(0.057) & \textbf{0.195\,(1.54)} \\
M08 & C+T+D+M & \textbf{0.560\,(6.8)} & \textbf{0.533\,(6.6)} & 0.446\,(0.057) & 0.188\,(1.57) & \textbf{0.612\,(6.4)} & \textbf{0.576\,(6.3)} & 0.391\,(0.058) & 0.147\,(1.58) \\
\midrule
\multicolumn{10}{l}{\textit{Panel\,C: 12-month window --- raw-lag memory}} \\
\midrule
M01 & C & 0.338\,(8.3) & 0.328\,(7.9) & 0.311\,(0.064) & 0.102\,(1.67) & 0.412\,(7.8) & 0.380\,(7.6) & 0.297\,(0.062) & 0.067\,(1.69) \\
M02 & C+D & 0.519\,(7.1) & 0.483\,(7.0) & 0.339\,(0.062) & 0.146\,(1.61) & 0.582\,(6.6) & 0.520\,(6.7) & 0.303\,(0.061) & 0.081\,(1.64) \\
M03 & C+T & 0.501\,(7.2) & 0.497\,(6.8) & 0.459\,(0.057) & 0.189\,(1.57) & 0.557\,(6.9) & 0.554\,(6.5) & 0.408\,(0.057) & 0.160\,(1.58) \\
M04 & C+T+D & 0.555\,(6.8) & 0.530\,(6.6) & 0.450\,(0.057) & 0.198\,(1.56) & 0.604\,(6.5) & 0.570\,(6.4) & 0.391\,(0.058) & 0.160\,(1.57) \\
M05 & C+M & \textbf{0.818\,(4.4)} & \textbf{0.841\,(3.9)} & \textbf{0.707\,(0.042)} & \textbf{0.666\,(1.05)} & \textbf{0.800\,(4.5)} & \textbf{0.824\,(4.1)} & \textbf{0.691\,(0.041)} & \textbf{0.622\,(1.09)} \\
M06 & C+D+M & 0.792\,(4.8) & 0.793\,(4.6) & 0.635\,(0.048) & 0.531\,(1.26) & 0.790\,(4.8) & 0.786\,(4.6) & 0.629\,(0.046) & 0.457\,(1.31) \\
M07 & C+T+M & 0.802\,(4.7) & 0.818\,(4.2) & 0.657\,(0.045) & 0.628\,(1.16) & 0.791\,(4.7) & 0.810\,(4.3) & 0.629\,(0.045) & 0.582\,(1.19) \\
M08 & C+T+D+M & 0.774\,(5.0) & 0.777\,(4.7) & 0.616\,(0.049) & 0.519\,(1.29) & 0.774\,(5.0) & 0.775\,(4.7) & 0.581\,(0.048) & 0.453\,(1.33) \\
\midrule
\multicolumn{10}{l}{\textit{Panel\,D: 24-month window --- raw-lag memory}} \\
\midrule
M01 & C & 0.344\,(8.3) & 0.326\,(8.0) & 0.318\,(0.063) & 0.102\,(1.67) & 0.414\,(7.8) & 0.362\,(7.8) & 0.305\,(0.062) & 0.069\,(1.69) \\
M02 & C+D & 0.518\,(7.1) & 0.472\,(7.0) & 0.325\,(0.063) & 0.118\,(1.64) & 0.584\,(6.6) & 0.510\,(6.8) & 0.281\,(0.062) & 0.067\,(1.66) \\
M03 & C+T & 0.502\,(7.2) & 0.490\,(6.9) & 0.461\,(0.056) & 0.193\,(1.57) & 0.566\,(6.8) & 0.541\,(6.6) & 0.413\,(0.057) & 0.167\,(1.57) \\
M04 & C+T+D & 0.559\,(6.8) & 0.533\,(6.6) & 0.441\,(0.058) & 0.205\,(1.56) & 0.611\,(6.4) & 0.572\,(6.4) & 0.384\,(0.058) & 0.168\,(1.56) \\
M05 & C+M & \textbf{0.838\,(4.1)} & \textbf{0.856\,(3.7)} & \textbf{0.734\,(0.040)} & \textbf{0.677\,(1.03)} & \textbf{0.818\,(4.3)} & \textbf{0.838\,(3.9)} & \textbf{0.713\,(0.039)} & \textbf{0.637\,(1.07)} \\
M06 & C+D+M & 0.822\,(4.5) & 0.821\,(4.2) & 0.681\,(0.045) & 0.593\,(1.19) & 0.808\,(4.6) & 0.807\,(4.4) & 0.671\,(0.044) & 0.533\,(1.23) \\
M07 & C+T+M & 0.834\,(4.3) & 0.841\,(3.9) & 0.711\,(0.042) & 0.640\,(1.11) & 0.817\,(4.4) & 0.827\,(4.1) & 0.686\,(0.042) & 0.590\,(1.16) \\
M08 & C+T+D+M & 0.811\,(4.6) & 0.809\,(4.4) & 0.663\,(0.046) & 0.587\,(1.21) & 0.798\,(4.7) & 0.798\,(4.5) & 0.634\,(0.046) & 0.528\,(1.25) \\
\bottomrule
\end{tabular}
\end{sidewaystable}


\end{document}